\documentclass[a4paper,12pt,twoside]{ctexart}
\usepackage[dvipsnames]{xcolor}
\usepackage{geometry}
\usepackage{tcolorbox}
\usepackage{tikz}
\usepackage{amssymb}
\usepackage{graphicx}

\definecolor{inkblack}{rgb}{0.08,0.08,0.08}
\geometry{a4paper, left=3.18cm, right=3.18cm, top=2.54cm, bottom=2.54cm}

\usepackage{fancyhdr}
\pagestyle{fancy}
\fancyhf{}
\renewcommand{\headrulewidth}{0pt}
\fancyfoot[RO]{\makebox[0pt][l]{\hspace*{28.2pt}\raisebox{2.81pt}{\rmfamily\fontsize{10.5bp}{10.5bp}\selectfont\thepage}}}
\fancyfoot[LE]{\makebox[0pt][r]{\hspace*{33.45pt}\raisebox{2.81pt}{\rmfamily\fontsize{10.5bp}{10.5bp}\selectfont\thepage}}}

\newCJKfontfamily\fzdbs[Path=./]{FZDBS.ttf}
\newCJKfontfamily\fzssk[Path=./]{FZSSK.ttf}
\AtBeginDocument{\color{inkblack}}

\newcommand{\lawtitle}[2]{
    \vspace*{0.507cm}
    \begin{center}
        {\fzdbs\fontsize{22bp}{22bp}\selectfont #1}\par
        \vspace{10.37bp}
        {\fzdbs\fontsize{18bp}{18bp}\selectfont #2}\par
    \end{center}
    \vspace{18.4bp}
}

\newtcolorbox{notesection}[1][]{
    width=467.72bp, sharp corners, colback=white, colframe=inkblack, boxrule=0.5pt,
    fontupper=\color{inkblack}\fontsize{10.5bp}{10.5bp}\selectfont,
    enlarge left by=-19.28bp, before=\vspace*{-13.39bp}\noindent, after=\par,
    boxsep=0pt, left=4bp, top=4.69bp, bottom=4.69bp, right=0pt, #1
}

\newcommand{\lawpar}{\par\vspace*{-4.86bp}}

\newcommand{\subtitlesection}[1]{%
    \begin{tcolorbox}[
        width=467.72bp, sharp corners, colback=white, colframe=inkblack, boxrule=0.5pt,
        boxsep=0pt, left=0pt, right=0pt, top=4bp, bottom=4bp, halign=center,
        fontupper=\fzdbs\fontsize{15bp}{15bp}\selectfont\color{inkblack},
        enlarge left by=-19.28bp,
        before={\nointerlineskip\vspace{-0.5pt}\noindent}, after=\vspace{0pt}
    ]
    #1
    \end{tcolorbox}%
}

\newcommand{\pairsection}[2]{%
    \begin{tcolorbox}[
        width=467.72bp, sharp corners, colback=white, colframe=inkblack, boxrule=0.5pt,
        enlarge left by=-19.28bp,
        before={\nointerlineskip\vspace{-0.5pt}\noindent}, after=\par,
        boxsep=0pt, left=0pt, right=0pt, top=0pt, bottom=0pt
    ]
    \begin{minipage}[c]{112.9bp}%
        \centering\songti\fontsize{10.5bp}{10.5bp}\selectfont #1%
    \end{minipage}%
    \hspace{0pt}%
    \vrule width 0.5pt%
    \hspace{4bp}%
    \begin{minipage}[c]{346.32bp}%
        \vspace{4bp}%
        \songti\mdseries\fontsize{10.5bp}{17bp}\selectfont%
        \baselineskip=17bp%
        \setlength{\parindent}{0pt}%
        \color{inkblack}%
        #2%
        \vspace{4bp}%
    \end{minipage}%
    \end{tcolorbox}%
}

\begin{document}

\lawtitle{民事起诉状}{（海上、通海水域人身损害责任纠纷）}

\begin{notesection}
\hfuzz=50pt
\color{inkblack}\heiti\fontsize{10.5bp}{16bp}\selectfont 说明：\par
\vspace*{-0.2bp}
\songti\mdseries\fontsize{10.5bp}{16bp}\selectfont
\setlength{\parindent}{20.7bp}
\sloppy
\setlength{\parskip}{0pt}
\vspace*{-5.1bp}为了方便您更好地参加诉讼，保护您的合法权利，请填写本表。
\lawpar
1.起诉时需向人民法院提交证明您身份的材料，如身份证复印件、营业执照复印件等。
\lawpar
2.本表所列内容是您提起诉讼以及人民法院查明案件事实所需，请务必如实填写。
\lawpar
3.本表有些内容可能与您的案件无关，您认为与案件无关的项目可以填"无"或不填；对于本\\[-4.86bp]
表中勾选项可以在对应项打"√"；您认为另有重要内容需要列明的，可以另附页填写。
\lawpar
4.本表 word 电子版填写时，相关栏目可复制粘贴或扩容，但不得改变要素内容、格式设置。例\\[-4.86bp]
如，多原告、多被告或多委托诉讼代理人等情况，可根据实际情况复制粘贴；需填写文字较多时，\\[-4.86bp]
可根据实际对栏目进行扩容等。
\lawpar
$\bigstar$特别提示$\bigstar$\\[-4.86bp]
\hspace*{20.7bp}诉讼参加人应遵守诚信原则如实认真填写表格。\\[-4.86bp]
\hspace*{20.7bp}如果诉讼参加人违反有关规定，虚假诉讼、恶意诉讼、滥用诉权，人民法院将视违法情形依法\\[-4.86bp]
追究责任。
\end{notesection}

\subtitlesection{当事人信息}

\pairsection{原告\\（自然人）}{
姓名：\\
性别：男$\square$\hspace{1em}女$\square$\\
出生日期：\hspace{2em}年\hspace{2em}月\hspace{2em}日\hspace{4em}民族：\\
工作单位：\hspace{6em}职务：\hspace{4em}联系电话：\\
住所地（户籍所在地）：\\
经常居住地：\\
证件类型：\\
证件号码：
}

\pairsection{原告\\（法人、非法人组织）}{
名称：\\
住所地（主要办事机构所在地）：\\
注册地 / 登记地：\\
法定代表人 / 负责人：\hspace{4em}职务：\hspace{4em}联系电话：\\
统一社会信用代码：\\
类型：有限责任公司$\square$\hspace{1em}股份有限公司$\square$\hspace{1em}上市公司$\square$\\
\hspace*{36bp}其他企业法人$\square$\hspace{1em}事业单位$\square$\hspace{1em}社会团体$\square$\hspace{1em}基金会$\square$\\
\hspace*{36bp}社会服务机构$\square$\hspace{1em}机关法人$\square$\hspace{1em}农村集体经济组织法人$\square$\\
\hspace*{36bp}城镇农村的合作经济组织法人$\square$\hspace{1em}基层群众性自治组织法人$\square$\\
\hspace*{36bp}个人独资企业$\square$\hspace{1em}合伙企业$\square$\hspace{1em}不具有法人资格的专业服务机构$\square$\\
所有制性质：国有$\square$（控股$\square$\hspace{1em}参股$\square$）\hspace{1em}民营$\square$\hspace{1em}其他\rule{60bp}{0.5pt}\\
外资情况：有$\square$\\
\hspace*{52bp}中外合资经营企业$\square$\hspace{1em}中外合作经营企业$\square$\\
\hspace*{52bp}外商独资企业$\square$\hspace{1em}其他形式\rule{60bp}{0.5pt}\\
\hspace*{52bp}无$\square$
}

\newpage

\pairsection{委托诉讼代理人}{
有$\square$\\
\hspace*{21bp}姓名：\\
\hspace*{21bp}\makebox[0.33\linewidth][l]{单位：}\makebox[0.33\linewidth][l]{职务：}\makebox[0.33\linewidth][l]{联系电话：}\\
\hspace*{21bp}代理权限：一般授权$\square$\hspace{1em}特别授权$\square$\hspace{1em}\rule{80bp}{0.5pt}\\
无$\square$
}

\pairsection{被告\\（自然人）}{
姓名：\\
性别：男$\square$\hspace{1em}女$\square$\\
出生日期：\hspace{2em}年\hspace{2em}月\hspace{2em}日\hspace{4em}民族：\\
工作单位：\hspace{6em}职务：\hspace{4em}联系电话：\\
住所地（户籍所在地）：\\
经常居住地：\\
证件类型：\\
证件号码：
}

\pairsection{被告\\（法人、非法人组织）}{
名称：\\
住所地（主要办事机构所在地）：\\
注册地 / 登记地：\\
法定代表人 / 负责人：\hspace{4em}职务：\hspace{4em}联系电话：\\
统一社会信用代码：\\
类型：有限责任公司$\square$\hspace{1em}股份有限公司$\square$\hspace{1em}上市公司$\square$\\
\hspace*{36bp}其他企业法人$\square$\hspace{1em}事业单位$\square$\hspace{1em}社会团体$\square$\hspace{1em}基金会$\square$\\
\hspace*{36bp}社会服务机构$\square$\hspace{1em}机关法人$\square$\hspace{1em}农村集体经济组织法人$\square$\\
\hspace*{36bp}城镇农村的合作经济组织法人$\square$\hspace{1em}基层群众性自治组织法人$\square$\\
\hspace*{36bp}个人独资企业$\square$\hspace{1em}合伙企业$\square$\hspace{1em}不具有法人资格的专业服务机构$\square$\\
所有制性质：国有$\square$（控股$\square$\hspace{1em}参股$\square$）\hspace{1em}民营$\square$\hspace{1em}其他\rule{60bp}{0.5pt}\\
外资情况：有$\square$\\
\hspace*{52bp}中外合资经营企业$\square$\hspace{1em}中外合作经营企业$\square$\\
\hspace*{52bp}外商独资企业$\square$\hspace{1em}其他形式\rule{60bp}{0.5pt}\\
\hspace*{52bp}无$\square$
}

\pairsection{第三人\\（自然人）}{
姓名：\\
性别：男$\square$\hspace{1em}女$\square$\\
出生日期：\hspace{2em}年\hspace{2em}月\hspace{2em}日\hspace{4em}民族：\\
工作单位：\hspace{6em}职务：\hspace{4em}联系电话：\\
住所地（户籍所在地）：\\
经常居住地：\\
证件类型：\\
证件号码：
}

\pairsection{第三人\\（法人、非法人组织）}{
名称：\\
住所地（主要办事机构所在地）：\\
注册地 / 登记地：\\
法定代表人 / 负责人：\hspace{4em}职务：\hspace{4em}联系电话：\\
统一社会信用代码：\\
类型：有限责任公司$\square$\hspace{1em}股份有限公司$\square$\hspace{1em}上市公司$\square$\\
\hspace*{36bp}其他企业法人$\square$\hspace{1em}事业单位$\square$\hspace{1em}社会团体$\square$\hspace{1em}基金会$\square$\\
\hspace*{36bp}社会服务机构$\square$\hspace{1em}机关法人$\square$\hspace{1em}农村集体经济组织法人$\square$\\
\hspace*{36bp}城镇农村的合作经济组织法人$\square$\hspace{1em}基层群众性自治组织法人$\square$\\
\hspace*{36bp}个人独资企业$\square$\hspace{1em}合伙企业$\square$\hspace{1em}不具有法人资格的专业服务机构$\square$\\
所有制性质：国有$\square$（控股$\square$\hspace{1em}参股$\square$）\hspace{1em}民营$\square$\hspace{1em}其他\rule{60bp}{0.5pt}\\
外资情况：有$\square$\\
\hspace*{52bp}中外合资经营企业$\square$\hspace{1em}中外合作经营企业$\square$\\
\hspace*{52bp}外商独资企业$\square$\hspace{1em}其他形式\rule{60bp}{0.5pt}\\
\hspace*{52bp}无$\square$
}

\newpage

\subtitlesection{诉讼请求}

\begin{tcolorbox}[
    width=467.72bp, sharp corners, colback=white, colframe=inkblack, boxrule=0.5pt,
    enlarge left by=-19.28bp,
    before={\nointerlineskip\vspace{-0.5pt}\noindent}, after=\par,
    boxsep=0pt, left=4bp, right=0pt, top=4bp, bottom=4bp
]
\songti\mdseries\fontsize{10.5bp}{17bp}\selectfont%
\baselineskip=17bp\parskip=0pt\lineskiplimit=-\maxdimen%
\setlength{\parindent}{0pt}\color{inkblack}%
（可完整表述诉讼请求；为方便、准确梳理要点，相关内容请在下方要素式表格中填写）\\[51bp]
\mbox{}
\end{tcolorbox}

\pairsection{1. 医疗费}{
\hspace{4em}年\hspace{2em}月\hspace{2em}日至\hspace{4em}年\hspace{2em}月\hspace{4em}日期间在\hspace{6em}医院住院（门诊）\\
治疗，累计支付医疗费\hspace{4em}元\\
医疗费发票、医疗费清单、病历资料：有$\square$\hspace{1em}无$\square$
}

\pairsection{2. 护理费}{
住院护理\hspace{2em}天支付护理费\hspace{4em}元（或护理人员发生误工费\hspace{4em}元），或遵医\\
嘱短期护理发生护理费\hspace{4em}元\\
住院证明、医嘱等：有$\square$\hspace{1em}无$\square$
}

\pairsection{3. 营养费}{
营养费\hspace{4em}元\\
病历资料：有$\square$\hspace{1em}无$\square$
}

\pairsection{4. 住院伙食补助费}{
住院伙食补助费\hspace{4em}元\\
病历资料：有$\square$\hspace{1em}无$\square$
}

\pairsection{5. 误工费}{
\hspace{4em}年\hspace{2em}月\hspace{2em}日至\hspace{4em}年\hspace{2em}月\hspace{2em}日误工费\hspace{4em}元
}

\pairsection{6. 交通费}{
交通费\hspace{4em}元\\
交通费凭证：有$\square$\hspace{1em}无$\square$
}

\pairsection{7. 残疾赔偿金（被扶养\\人生活费计入）}{
残疾赔偿金\hspace{4em}元（含被扶养人生活费\hspace{4em}元）
}

\pairsection{8. 残疾辅助器具费}{
残疾辅助器具费\hspace{4em}元
}

\pairsection{9. 死亡赔偿金（被扶养\\人生活费计入）、丧葬费}{
死亡赔偿金\hspace{4em}元（含被扶养人生活费\hspace{4em}元），丧葬费\hspace{4em}元
}

\pairsection{10. 精神损害抚慰金}{
精神损害抚慰金\hspace{4em}元
}

\pairsection{11. 是否主张诉讼费用}{
是$\square$\\
否$\square$
}

\pairsection{12. 其他费用}{
主张\hspace{8em}费用\hspace{4em}元
}

\pairsection{13. 赔偿总额}{
\hspace{4em}元（计至\hspace{4em}年\hspace{2em}月\hspace{2em}日）
}

\pairsection{14. 是否主张船舶优先权}{
是$\square$\hspace{1em}内容：\\
否$\square$
}

\subtitlesection{约定管辖和诉前保全}

\pairsection{1. 有无仲裁、法院管辖\\约定}{
有$\square$\hspace{4em}合同条款及内容：\\
无$\square$
}

\pairsection{2. 是否已经诉前保全}{
是$\square$\hspace{1em}保全法院：\hspace{4em}保全时间：\\
\hspace*{21bp}保全案号：\\
否$\square$\\
（如申请诉讼保全，请另行提交诉讼保全申请及相关材料）
}

\newpage

\subtitlesection{涉外及涉港澳台}

\pairsection{1. 是否涉外}{
是$\square$\hspace{1em}涉及国家：\\
否$\square$
}

\pairsection{2. 是否涉港澳台}{
是$\square$\hspace{1em}涉港$\square$\hspace{1em}涉澳$\square$\hspace{1em}涉台$\square$\\
否$\square$
}

\subtitlesection{事实与理由}

\begin{tcolorbox}[
    width=467.72bp, sharp corners, colback=white, colframe=inkblack, boxrule=0.5pt,
    enlarge left by=-19.28bp,
    before={\nointerlineskip\vspace{-0.5pt}\noindent}, after=\par,
    boxsep=0pt, left=4bp, right=0pt, top=4bp, bottom=4bp
]
\songti\mdseries\fontsize{10.5bp}{17bp}\selectfont%
\baselineskip=17bp\parskip=0pt\lineskiplimit=-\maxdimen%
\setlength{\parindent}{0pt}\color{inkblack}%
（可完整表述纠纷涉及的事实与理由；为方便、准确梳理要点，相关内容请在下方要素式表格中填写）\\[51bp]
\mbox{}
\end{tcolorbox}

\pairsection{1. 请求依据}{
合同约定：\\
法律规定：
}

\pairsection{2. 船舶情况}{
船名：\hspace{8em}船舶所有人：
}

\pairsection{3. 劳务关系情况}{
雇主姓名：\hspace{8em}工资标准：\\
上船时间：\hspace{8em}下船时间：
}

\pairsection{4. 事故情况}{
死亡$\square$\hspace{1em}失踪$\square$\hspace{1em}无伤残等级$\square$\hspace{1em}有伤残等级$\square$\hspace{2em}级\\
事故调查机关：\\
事故原因：船上劳务$\square$\hspace{1em}船上劳务相关活动$\square$\hspace{1em}其他\rule{60bp}{0.5pt}
}

\pairsection{5. 船舶、船员投保情况}{
\mbox{}
}

\pairsection{6. 其他情况}{
\mbox{}
}

\pairsection{7. 证据清单（可另附\\页）}{
\mbox{}
}

\subtitlesection{对纠纷解决方式的意愿}

\pairsection{是否了解调解作为非诉\\讼纠纷解决方式，能\\及时、高效、低成本、\\不伤和气地解决纠纷}{
了解$\square$\hspace{1em}不了解$\square$
}

\pairsection{是否了解先行调解解\\决纠纷的好处}{
1. 立案后选择先行调解的，可以很快启动调解程序。如不同意调解，法院将\\
依程序开庭审理案件，但可能需要经过较长一段时间的排期等待，且审理、\\
执行周期相对较长。\\
了解$\square$\hspace{1em}不了解$\square$\\
2. 选择先行调解，调解成功且自动履行的免交诉讼费用，申请司法确认的不\\
交纳诉讼费用，要求出具调解书的减半交纳诉讼费用。\\
了解$\square$\hspace{1em}不了解$\square$\\
3. 首次调解不成功，但仍有继续调解意愿的，可以选择更换调解组织和调解\\
员再进行调解。调解无法达成一致意见的，法院将依程序排期开庭。\\
了解$\square$\hspace{1em}不了解$\square$\\
4. 依照法律规定，调解具有保密性要求，调解过程不公开，调解协议未经当\\
事人同意不得公开。\\
了解$\square$\hspace{1em}不了解$\square$\\
5. 调解达成的协议具有法律效力，可以依照法律规定申请司法确认，具有强\\
制执行效力。\\
了解$\square$\hspace{1em}不了解$\square$
}

\pairsection{是否考虑先行调解}{
是$\square$\\
否$\square$\\
暂不确定，想要了解更多内容$\square$
}

\vspace{10bp}
\noindent\hspace*{248bp}{\fzdbs\fontsize{15bp}{22bp}\selectfont 具状人（签字、盖章）：}\par
\noindent\hspace*{248bp}{\fzdbs\fontsize{15bp}{22bp}\selectfont 日期：}

\end{document}
